\section{Angular determinant and electroweak-sector parameters}
\label{rm:ch:ew}

\subsection{Constitutive and Determinantal Functions of the Angular Operator}

Recall the two-sheet operator of the \cref{rm:ch:vacuum},
\begin{equation}
\label{rm:eq:ew-angular-operator}
T=q_+P_+ +q_-P_-,
\qquad
q_-=\ee^{-(A_{\rm rad}-C^*_{\rm rad})},
\qquad
q_+=\ee^{-(A_{\rm rad}+C^*_{\rm rad})}.
\end{equation}
The electromagnetic vacuum applies the function
\begin{equation}
\label{rm:eq:ew-vacuum-reader}
F(T)^2,
\qquad
F(z)=\frac{1-z^{90}}{1-z^{120}},
\end{equation}
while the electroweak angular sector applies the determinant
\begin{equation}
\label{rm:eq:ew-determinant-reader}
\boxed{
\det T=q_+q_-=\ee^{-2A_{\rm rad}}=:D_A.}
\end{equation}
They are two sibling functional realizations of the common antecedent
\((A,C^*,T)\). The causal structure generates both from \(T\). The
subsequent reconstruction demonstrated in
\cref{rm:thm:curve-F-bijection,rm:thm:curve-vacuum-tomography} permits recovery of
\[
 \det T
 =\det\!\left(F^{-1}\!\left(\sqrt{F(T)^2}\right)\right),
\]
without turning that functional inversion into the generating arrow of the
electroweak sector. Structural unity resides in \(T\), and the reconstruction
makes explicit the compatibility among permeability, permittivity, and
\(\theta_W\).

Under the sheet involution \(C^*\mapsto-C^*\), the eigenvalues
\(q_+,q_-\) are exchanged. Therefore \(D_A\) and the constitutive speed
are even, while \((\widehat\mu,\widehat\varepsilon)\) is exchanged and
\(\widehat Z\) is inverted. The following electroweak relation uses the
even invariant.

\subsection{Exact relation between the \texorpdfstring{\(W\)}{W} and \texorpdfstring{\(Z\)}{Z} sectors}

The angular coordinate of duration--perduration is
\begin{equation}
\label{rm:eq:ew-angular-character}
\chi_{\rm dur}(n_A,s_C)
=\exp\!\left(n_AA_{\rm rad}
+\frac{s_C}{6}C^*_{\rm rad}\right).
\end{equation}
This coordinate is not introduced as an isolated label. Its origin is
\begin{equation}
\label{rm:eq:ew-signature-provenance}
\begin{aligned}
\APP&\longrightarrow\TRIT\longrightarrow\TPK
\longrightarrow\text{surviving extension}
\longrightarrow\text{route}\\
&\longrightarrow\text{genealogical record}
\longrightarrow(n_A,s_C,\ldots)
\longrightarrow\chi_{\rm dur}.
\end{aligned}
\end{equation}
The signatures of \(W\), \(Z\) and \(H\) are typed quantities obtained through
that chain; this chapter adopts them without selecting them from
their terminal masses. In the established angular realization, the
charged and neutral sectors satisfy
\begin{equation}
\label{rm:eq:ew-WZ-signatures}
(n_A,s_C)_W=(94,-1),
\qquad
(n_A,s_C)_Z=(95,-1).
\end{equation}
This assertion refers to the common angular component. The
additional refinement coordinates of the General Mass Law are not
eliminated unless the physical transducer proves their cancellation.

\begin{theorem}[Weinberg defect in the physical-mass scheme of the angular chart]
\label{rm:thm:ew-weinberg}
If \(W\) and \(Z\) are evaluated on the same scale section and share
the refinement that accompanies \cref{rm:eq:ew-WZ-signatures}, then
\begin{equation}
\label{rm:eq:ew-WZ-ratio}
\boxed{
\frac{m_W^{\angle}}{m_Z^{\angle}}
=\ee^{-A_{\rm rad}}=\sqrt{D_A},
\qquad
\left(\frac{m_W^{\angle}}{m_Z^{\angle}}\right)^2=D_A.}
\end{equation}
In the physical-mass scheme,
\begin{equation}
\label{rm:eq:ew-sin2thetaW}
\boxed{
\sin^2\theta_W^{\rm OS,HMT}=1-D_A
=0.2248708807528435698111159035\ldots .}
\end{equation}
Adem\'as,
\begin{equation}
\label{rm:eq:ew-WZ-terminal-ratio}
\frac{m_W^{\angle}}{m_Z^{\angle}}
=0.8804141748331613670128885459\ldots .
\end{equation}
\end{theorem}

\begin{proof}
By \cref{rm:eq:ew-angular-character,rm:eq:ew-WZ-signatures},
\[
\frac{\chi_{\rm dur}(94,-1)}{\chi_{\rm dur}(95,-1)}
=\exp(-A_{\rm rad});
\]
the coordinate \(C^*\) cancels because both sectors share
\(s_C=-1\). Squaring and using
\cref{rm:eq:ew-determinant-reader} yields \(D_A\). The definition
of physical masses \(s_W^2=1-m_W^2/m_Z^2\) yields
\cref{rm:eq:ew-sin2thetaW}.
\end{proof}

The principal content is the identity between two internal constructions: the determinant of the angular operator and the one-step difference in the realization \(W/Z\). Promotion to a physical gauge selector requires an intertwiner that carries the sheets of \(T\) to the neutral basis \((W^3,B)\) and demonstrates which refinements are preserved or cancel.

\subsection{Gauge couplings in rationalized normalization}

We define the dimensionless gauge coupling
\begin{equation}
\label{rm:eq:ew-egauge}
e_g:=\sqrt{4\pi\alpha_{\HMT}}
=0.3028221208717526427662442689\ldots .
\end{equation}
\(e_g\) must not be confused with a dimensional charge section
\(e_{\rm int}\): the former is a coupling in natural units; the
latter lives in the charge line of Dimensional Mechanics.

Let
\begin{equation}
\label{rm:eq:ew-sw-cw}
s_W^2=1-D_A,
\qquad
c_W^2=D_A.
\end{equation}
The convention electrodebil racionalizada to level of arbol gives
\begin{align}
g&=\frac{e_g}{s_W}
=2\sqrt{\frac{\pi\alpha_{\HMT}}{1-D_A}}
=0.6385883427334574283647003809\ldots,
\label{rm:eq:ew-g}\\
g'&=\frac{e_g}{c_W}
=2\sqrt{\frac{\pi\alpha_{\HMT}}{D_A}}
=0.3439541633108502429362319257\ldots,
\label{rm:eq:ew-gprime}\\
\frac{g'}{g}
&=\sqrt{D_A^{-1}-1}
=0.5386164141966093594996610933\ldots .
\label{rm:eq:ew-gprime-over-g}
\end{align}

\begin{proposition}[Tree-level custodial relation]
\label{rm:prop:ew-rho}
At the same interface,
\begin{equation}
\label{rm:eq:ew-rho}
\rho_{\rm EW}^{(0)}
:=\frac{(m_W^{\angle})^2}
{(m_Z^{\angle})^2c_W^2}=1.
\end{equation}
\end{proposition}

\begin{proof}
By \cref{rm:eq:ew-WZ-ratio,rm:eq:ew-sw-cw}, both the relative numerator and
\(c_W^2\) are \(D_A\).
\end{proof}

The equations \cref{rm:eq:ew-g,rm:eq:ew-gprime} are exact compositions once the
physical mass scheme, rationalization, and scale have been declared. They do not
assert that a change of renormalization scheme leaves their
numerical values unchanged.

\subsection{Fermi Constant in the Approximation at Tree Level}

In natural units, integration of the charged propagator at low
energy yields
\begin{equation}
\label{rm:eq:ew-GF-def}
\frac{G_F^{(0)}}{\sqrt2}=\frac{g^2}{8m_W^2}.
\end{equation}
Substituting \cref{rm:eq:ew-g} gives the composition
\begin{equation}
\label{rm:eq:ew-GF}
\boxed{
G_F^{(0)}
=\frac{\pi\alpha_{\HMT}}
{\sqrt2(1-D_A)m_W^2}.}
\end{equation}
By adopting the previously obtained mass
\begin{equation}
\label{rm:eq:ew-W-output}
m_W^{\HMT}=80.381592550221\ldots\ \mathrm{GeV},
\end{equation}
the terminal evaluation is
\begin{equation}
\label{rm:eq:ew-GF-value}
\boxed{
G_F^{(0)}
=1.1157162817659203186621784\times10^{-5}\ \mathrm{GeV}^{-2}.}
\end{equation}

The difference with respect to a renormalized Fermi constant is not used
to choose any coefficient. Under the convention
\begin{equation}
\label{rm:eq:ew-Delta-r}
G_F
=\frac{\pi\alpha_{\HMT}}
{\sqrt2(1-D_A)m_W^2(1-\Delta r)},
\end{equation}
the remaining obligation is to derive \(\Delta r\) from the internal sector of
loops, states, and scale. Opening the observed value of \(G_F\) to define
\(\Delta r\) would be a calibrated reconstruction, not a prediction.

\subsection{Higgs Coupling Oriented Dependence}

In the angular chart, the relevant routes are
\begin{equation}
\label{rm:eq:ew-HW-signatures}
(n_A,s_C)_H=(97,9),
\qquad
(n_A,s_C)_W=(94,-1).
\end{equation}
Therefore,
\begin{equation}
\label{rm:eq:ew-HW-ratio}
\boxed{
\frac{m_H^{\angle}}{m_W^{\angle}}
=\frac{\chi_{\rm dur}(97,9)}{\chi_{\rm dur}(94,-1)}
=\exp\!\left(3A_{\rm rad}+\frac53C^*_{\rm rad}\right)
=1.55872087212325548564\ldots .}
\end{equation}
Here, \(C^*\) does not cancel: the scalar sector preserves the orientation
that the determinant \(D_A=\ee^{-2A_{\rm rad}}\) does not distinguish.

Let us take the convention of the potential
\begin{equation}
\label{rm:eq:ew-Higgs-potential}
V(H)=-\mu_H^2H^\dagger H
+\lambda_H(H^\dagger H)^2,
\end{equation}
for which
\begin{equation}
\label{rm:eq:ew-Higgs-tree-relations}
m_H^2=2\lambda_Hv^2,
\qquad
m_W^2=\frac{g^2v^2}{4}.
\end{equation}

\begin{corollary}[Oriented Higgs self-coupling]
\label{rm:cor:ew-Higgs-lambda}
At the declared tree-level interface,
\begin{equation}
\label{rm:eq:ew-Higgs-lambda}
\boxed{
\lambda_H
=\frac{\pi\alpha_{\HMT}}{2(1-D_A)}
\exp\!\left(6A_{\rm rad}+\frac{10}{3}C^*_{\rm rad}\right)
=0.1238479115482466804662\ldots .}
\end{equation}
Adem\'as,
\begin{equation}
\label{rm:eq:ew-GF-Higgs}
G_F^{(0)}m_H^2=\sqrt2\,\lambda_H.
\end{equation}
\end{corollary}

\begin{proof}
Dividing the two equations in
\cref{rm:eq:ew-Higgs-tree-relations} gives
\(m_H^2/m_W^2=8\lambda_H/g^2\). With
\(g^2=4\pi\alpha_{\HMT}/(1-D_A)\) and the square of
\cref{rm:eq:ew-HW-ratio}, we obtain \cref{rm:eq:ew-Higgs-lambda}. On the other hand,
\(G_F^{(0)}=1/(\sqrt2v^2)\) and
\(m_H^2=2\lambda_Hv^2\) imply
\cref{rm:eq:ew-GF-Higgs}.
\end{proof}

The value of \(\lambda_H\) precedes the incorporation of radiative corrections and depends on the declared convention and scale. Its structural interest is that it separates two memories: the common gauge sector uses the even invariant \(D_A\), whereas the Higgs route explicitly retains \(C^*\).

\subsection{Determinantal closure of the electroweak sector: mixing, couplings, and constitutive spectrum}

The complete chain of this chapter is
\begin{equation}
\label{rm:eq:ew-causal-diagram}
\boxed{
\begin{aligned}
(A,C^*)&\longrightarrow T
\longrightarrow
\begin{cases}
F(T)^2&\longrightarrow
(\widehat\mu,\widehat\varepsilon,\widehat Z,\widehat c),\\
\det T=D_A&\longrightarrow
(s_W^2,c_W^2),
\end{cases}\\
(D_A,\alpha_{\HMT})&\longrightarrow(e_g,g,g'),\\
(D_A,\alpha_{\HMT},m_W)&\longrightarrow G_F^{(0)},\\
(A,C^*,D_A,\alpha_{\HMT})&\longrightarrow
\left(\frac{m_H}{m_W},\lambda_H\right).
\end{aligned}}
\end{equation}
The figure always appears at the end. Neither \(\sin^2\theta_W\), nor
\(G_F\), nor \(\lambda_H\) is introduced as a target value for selecting
\(A,C^*\) or the route signatures.

\begin{keyresult}
The same angular operator determines the vacuum through a rational function
of its sheets and the electroweak defect through its determinant. The
one-step difference between \(W\) and \(Z\) turns that determinant into
\(\sin^2\theta_W^{\rm OS}=1-D_A\); the Higgs route, in contrast, preserves
the counterangle and preserves the orientation eliminated by the determinant.
\end{keyresult}

\begin{falsifier}
The Weinberg relation is refuted in its domain if the angular
signatures \(W/Z\) do not differ only by \(n_A:94\to95\), if a refinement
coordinate does not cancel at the declared interface, or if
\((m_W^{\angle}/m_Z^{\angle})^2\neq D_A\). Gauge promotion is
blocked if no intertwiner with the basis \((W^3,B)\) exists. The figures
of \(g,g',G_F,\lambda_H\) change if the scheme or scale is altered without
recording it. Using observed \(G_F\) or \(m_H\) to select
\(\Delta r\), the signatures, or \(C^*\) introduces feedback from the
target value and invalidates the derivation. Conflating \(e_g\) with the dimensional
charge constitutes a type error.
\end{falsifier}

\begin{statusbox}
Angular determinant, quotient \(W/Z\), and complementary Weinberg defect:
\status{Exact internal theorem in the angular conditioned chart}. Couplings \(g,g'\),
\(\rho_{\rm EW}^{(0)}\), Fermi, and Higgs: \status{Exact Composition in the
Electroweak Tree-Level Interface}. The gauge-basis intertwiner and the
internal correction \(\Delta r\): \status{Open boundary with falsification criterion}. Fermi and Higgs are not independent constants selected without target values: they are exact compositions using the signatures, the quantity \(m_W\) and the declared tree-level convention. The subsequent conventional values are \status{Contrast metrological external}.
\end{statusbox}
